\DocumentMetadata{
  pdfversion=2.0,
  pdfstandard=ua-2,
  lang=en-US,
  tagging=on,
  tagging-setup={math/setup=mathml-SE,tabsorder=structure}
}
\documentclass[11pt,letterpaper]{article}
\usepackage[margin=0.68in,includefoot]{geometry}
\usepackage{newcomputermodern}
\usepackage{amsmath,amscd,mathtools}
\usepackage{tikz,adjustbox}
\providecommand{\square}{\mdlgwhtsquare}
\usepackage[hidelinks]{hyperref}
\hypersetup{
  pdftitle={Algebra PhD Exam, August 2018},
  pdfauthor={University of Florida Department of Mathematics},
  pdfsubject={Graduate mathematics examination},
  pdfdisplaydoctitle=true,
  bookmarksopen=true
}
\setcounter{secnumdepth}{0}
\setlength{\parindent}{0pt}
\setlength{\parskip}{0.28em}
\setlength{\emergencystretch}{3em}
\setlength{\leftmargini}{2.15em}
\setlength{\leftmarginii}{2.65em}
\setlength{\leftmarginiii}{3em}
\pagestyle{empty}
\raggedbottom
\allowdisplaybreaks
\NewTaggingSocketPlug{block/list/label}{examproblem}{%
  \tagstructbegin{tag=Lbl}\tagstructend%
  \tagstructbegin{tag=\LBody}%
  \tagstructbegin{tag=\ProblemHeadingTag,title-o={\ProblemHeadingText}}%
  \tagmcbegin{tag=\ProblemHeadingTag,actualtext={\ProblemHeadingText}}%
  #2%
  \tagmcend%
  \tagstructend%
}
\newenvironment{examproblems}[2]{%
  \def\ProblemHeadingTag{#2}%
  \AssignTaggingSocketPlug{block/list/label}{examproblem}%
  \begin{enumerate}%
  \setcounter{enumi}{#1}%
  \renewcommand{\labelenumi}{\textbf{\theenumi.}}%
  \setlength{\topsep}{0.35em}%
  \setlength{\partopsep}{0pt}%
  \setlength{\itemsep}{0.48em}%
  \setlength{\parsep}{0.18em}%
}{\end{enumerate}}
\newenvironment{examsubparts}{%
  \AssignTaggingSocketPlug{block/list/label}{default}%
  \begin{enumerate}%
  \setlength{\topsep}{0.18em}%
  \setlength{\partopsep}{0pt}%
  \setlength{\itemsep}{0.16em}%
  \setlength{\parsep}{0.1em}%
}{\end{enumerate}}
\newenvironment{examinstructions}{%
  \par\begingroup\itshape
}{%
  \par\endgroup\vspace{0.18em}
}
\makeatletter
\renewcommand\subsection{\@startsection{subsection}{2}{0pt}%
  {-1.25ex plus -.25ex minus -.1ex}%
  {0.55ex plus .1ex}%
  {\normalfont\large\bfseries}}
\renewcommand{\@maketitle}{%
  \newpage\null\vskip -1em%
  {\footnotesize\bfseries
    \href{https://www.ufl.edu/}{University of Florida}, \href{https://math.ufl.edu/}{Department of Mathematics}\par}%
  \vskip -0.5em%
  \tagstructbegin{tag=H1,title={Algebra PhD Exam, August 2018}}%
  \tagpdfparaOff%
  \tagmcbegin{tag=H1}%
    {\LARGE\bfseries\href{https://gma.math.ufl.edu/exams/algebra-phd/}{Algebra PhD Exam}, August 2018\par}%
  \tagmcend%
  \tagpdfparaOn%
  \tagstructend%
  \vskip 0.25em%
  \tagmcbegin{artifact}%
    \hrule height 1pt%
  \tagmcend%
  \vskip 1em%
}
\makeatother
\title{Algebra PhD Exam, August 2018}
\author{}
\date{}
\begin{document}
\maketitle
\thispagestyle{empty}
\begin{examinstructions}
Answer seven problems. Write your answers clearly in complete English sentences. You may quote results (within reason) as long as you state them clearly.
\end{examinstructions}
\begin{examproblems}{0}{H2}
\def\ProblemHeadingText{Problem 1.}
\item
Determine the Galois group of the polynomial
\[
f=x^5-2\in\mathbb{Q}[x].
\]
\def\ProblemHeadingText{Problem 2.}
\item
Let~\(F\) be a field of characteristic zero and~\(E\) a finite extension field of~\(F\). Prove that~\(E\) is a simple extension of~\(F\), i.e. there exists \(\alpha\in E\) such that \(E=F(\alpha)\).
\def\ProblemHeadingText{Problem 3.}
\item
Let~\(R\) be a ring with 1. Prove the equivalence of the following conditions.
\begin{examsubparts}
\item[\textnormal{(a)}]
Every unital left~\(R\)-module is projective.
\item[\textnormal{(b)}]
Every short exact sequence of unital left~\(R\)-modules splits.
\item[\textnormal{(c)}]
Every unital left~\(R\)-module is injective.
\end{examsubparts}
\def\ProblemHeadingText{Problem 4.}
\item
Prove that the additive group of the field of rational numbers is not a projective~\(\mathbb{Z}\)-module.
\def\ProblemHeadingText{Problem 5.}
\item
Let~\(R\) and~\(S\) be rings, \(A\) a right~\(R\)-module, \(B\) an \((R,S)\)-bimodule and~\(C\) a left~\(S\)-module.
\begin{examsubparts}
\item[\textnormal{(a)}]
Explain how \(B\otimes_S C\) is a left~\(R\)-module and \(A\otimes_R B\) is a right~\(S\)-module.
\item[\textnormal{(b)}]
Prove that we have an isomorphism (of abelian groups)
\[
A\otimes_R(B\otimes_S C)\cong(A\otimes_R B)\otimes_S C.
\]
\end{examsubparts}
\def\ProblemHeadingText{Problem 6.}
\item
Let~\(R\) be an integral domain and for each maximal ideal~\(M\), regard the localization \(R_M\) as a subring of the field of fractions~\(F\) of~\(R\). Show that \(\bigcap_M R_M=R\). (Hint: for \(x\in F\), consider the ideal \(D=\{a\in R\mid ax\in R\}\). If \(x\in R_M\) what does it say about~\(D\)?)
\def\ProblemHeadingText{Problem 7.}
\item
Let~\(R\) be a Dedekind domain and~\(I\) a nonzero ideal of~\(R\).
\begin{examsubparts}
\item[\textnormal{(a)}]
Prove that \(R/I\) is Artinian.
\item[\textnormal{(b)}]
Prove that if~\(R\) is itself Artinian, then it is a field.
\end{examsubparts}
\def\ProblemHeadingText{Problem 8.}
\item
Let \(F[[X]]\) denote the ring of formal power series over the field~\(F\).
\begin{examsubparts}
\item[\textnormal{(a)}]
Prove that \(a_0+a_1X+a_2X^2+\cdots\in F[[X]]\) is a unit if and only if \(a_0\ne0\).
\item[\textnormal{(b)}]
Prove that \(F[[X]]\) is a discrete valuation ring.
\end{examsubparts}
\def\ProblemHeadingText{Problem 9.}
\item
Let~\(R\) and~\(S\) be commutative rings with 1. Prove the Lying-Over Theorem: Let~\(S\) be an integral extension of~\(R\) and~\(P\) a prime ideal of~\(R\). Then there exists a prime ideal~\(Q\) of~\(S\) such that \(Q\cap R=P\).
\def\ProblemHeadingText{Problem 10.}
\item
Prove that the ring of \(n\times n\) matrices with entries in a field~\(F\) is a simple ring when \(n\ge1\). Describe explicitly one of its minimal left ideals.
\def\ProblemHeadingText{Problem 11.}
\item
Let \(\mathcal C\) be a category and let~\(\Lambda\) be a set. Suppose that for each \(\lambda\in\Lambda\) we are given an object \(X_\lambda\) in \(\mathcal C\).
\begin{examsubparts}
\item[\textnormal{(a)}]
Give the definition of a product of the collection of objects \(\{X_\lambda\}_{\lambda\in\Lambda}\).
\item[\textnormal{(b)}]
Prove that any two products of the \(\{X_\lambda\}_{\lambda\in\Lambda}\) are isomorphic.
\item[\textnormal{(c)}]
Prove that in the category of Groups a product exists for any collection of objects indexed by a set~\(\Lambda\).
\item[\textnormal{(d)}]
Prove that if, in (c), the groups in the collection happen to be abelian, then a product as in (c) is also a product in the category of Abelian groups.
\end{examsubparts}
\end{examproblems}
\end{document}
